Todos estos resultados se unen porque convierten el valor terminal en la expresión cuantitativa de una estructura previamente construida. La exactitud numérica acompaña a una concepción distinta de qué es una constante fundamental. La conexión nonádica no genera primero tres constantes para añadir después una cuarta: la prolongación \(w_6\to w_{12}\to w_{18}\to w_{24}\to w_{30}\to R_{36}\to G_9\) estabiliza un estado coinductivo común del que se publican correlacionadamente \(\pi\), \(e\), \(\varphi\) y \(\alpha\). Sus lectores distinguen clausura, propagación, autoescala y acoplamiento; la coordenada \(\alpha\) conserva, dentro de esa generación conjunta, su realización firmada y dodecafásica y sus dos vías internas de derivación.\par

En la descripción habitual, una constante suele presentarse como un número situado delante de una ecuación: una cantidad que la teoría necesita, pero cuyo valor debe recibirse desde fuera. En Holografía Modular Triádica ocurre lo contrario. La constante aparece al final, cuando un mismo estado ha sido recorrido por varios lectores y se comprueba qué relación ha sobrevivido a todos ellos.\par

Por eso las constantes del documento son magnitudes dotadas de genealogía: restos invariantes de una transformación. Cada una dice qué se ha conservado cuando algo ha cambiado:\par

\begin{itemize}[leftmargin=*,itemsep=0.35em]
\item una hoja se ha invertido;
\item un ciclo ha regresado a su fase inicial;
\item una memoria ha atravesado la costura;
\item una descripción interior ha sido publicada en el borde;
\item una magnitud adimensional ha recibido una realización física;
\item un sistema ha cambiado de escala conservando íntegramente su genealogía.
\end{itemize}

La Mecánica Dimensional añade después las rectas de magnitud y las unidades. Pero la estructura que selecciona la constante es anterior. El valor real aparece al final porque \(\mathbb R\) funciona aquí como evaluación terminal: es la sombra cuantitativa de una construcción que ya estaba completa antes de convertirse en decimal.\par

Con esa visión, los nueve desarrollos considerados forman una única narración sobre orientación, memoria, geometría, información y escala.\par
